%!TEX root = ../chapter04.tex 
%----------------------------------------------------------------------------------------
%	
%----------------------------------------------------------------------------------------


%----------------------------------------------------------------------------------------
%	 
%----------------------------------------------------------------------------------------

\newpage
\loadgeometry{widelayout}  
\pagestyle{scrheadings} % Print headers again  
\KOMAoptions{
	headwidth=textwithmarginpar,
	headsepline=true,
	headsepline= 0.5pt,
	footwidth=textwithmarginpar,
}   
\onehalfspacing   


\section{TD Implications et raisonnement par contraposée}
$P$ et $Q$ sont des affirmations dont les valeurs de vérité sont soit VRAI soit FAUX. 
Une implication est un énoncé de la forme \og Si $P$ alors $Q$\fg{} ou \og $P$ implique $Q$\fg{} (en abrégé $P\Rightarrow Q$).
\begin{tBox} 
	\og  $P$ implique $Q$ \fg{} est un énoncé FAUX quand \og $P$ VRAI et $Q$ FAUX \fg{}. 	\\
	L'implication est VRAI dans les autres cas. 
\end{tBox}
\begin{example}Donner les valeurs de vérité des implications suivantes :\hfill 
\end{example}
\QCM[VF,Alterne,Stretch=1.1, Largeur=0.75cm]{%
	{Si \uline{12 est un multiple de 9}, alors \uline{12 est un multiple de 3}.}&1,%
	{Si \uline{10 est un multiple de 9}, alors \uline{10 est un multiple de 3}.}&1,%
	{Si \uline{$x$ est un multiple de 9}, alors \uline{$x$ est un multiple de 3}.}&1,%
	{Si \uline{18 est un multiple de 3}, alors \uline{18 est un multiple de 9}.}&1,%
	{Si \uline{12 est un multiple de 3}, alors \uline{12 est un multiple de 9}.}&2,%
	{Si \uline{10 est un multiple de 3}, alors \uline{10 est un multiple de 9}.}&1,%
	{Si \uline{$x$ est un multiple de 3}, alors \uline{$x$ est un multiple de 9}.}&2,%
}
\begin{tBox}[frametitle=À retenir]
	\noindent\parbox{0.31\linewidth}{\centering
		Sa contraposée
		
		Si $\operatorname{non}Q$ alors $\operatorname{non}P$
		
	}\hfill \parbox{0.31\linewidth}{\centering
		L'implication
		
		Si $P$ alors $Q$
		
	}\hfill \parbox{0.31\linewidth}{\centering
		Sa réciproque
		
		Si $Q$ alors $P$ 
	}
\end{tBox}	 
\begin{example}  	Pour chacune des implications suivantes rédiger la réciproque et la contraposée et préciser si elles sont vraies ou fausses.  
\end{example}
\noindent\QCM[VF,Alterne,Stretch=1.1, Largeur=0.75cm]{%
	{\textbf{Implication} Si \uline{ c'est un ours polaire}, alors \uline{l'animal sait nager}.}&1,%
	{\textbf{Contraposée} }&1,% 
	{\textbf{Réciproque} }&2,% 
}	 

\noindent\QCM[VF,Alterne,Stretch=1.1, Largeur=0.75cm]{%
	{\textbf{Implication} Si \uline{le triangle $ABC$ est rectangle en $A$}, alors \uline{$AB^2+AC^2=BC^2$}.}&1,%
	{\textbf{Contraposée } }&1,% 
	{\textbf{Réciproque } }&2,% 
}


\begin{exercise}[À vous]\label{chap03:exercice:logique01} Mêmes consignes
	
	\noindent\QCM[VF,Alterne,Stretch=1.1, Largeur=0.75cm]{%
		{\textbf{Implication} Si \uline{l'animal est une poule}, alors \uline{l'animal pond des oeufs}.}&1,%
		{\textbf{Contraposée } }&1,% 
		{\textbf{Réciproque } }&2,% 
	}
	
	
	\noindent
	\QCM[VF,Alterne,Stretch=1.1, Largeur=0.75cm]{%
		{\textbf{Implication} Si \uline{c’est un carré} alors \uline{le quadrilatère a 4 angles droits}.}&1,%
		{\textbf{Contraposée } }&1,% 
		{\textbf{Réciproque } }&2,% 
	}	
	
	\noindent
	\QCM[VF,Alterne,Stretch=1.1, Largeur=0.75cm]{%
		{\textbf{Implication} Si \uline{deux droites du plan se coupent} alors \uline{elles sont perpendiculaires}.}&1,%
		{\textbf{Contraposée } \textcolor{gray!15}{Si les deux droites ne sont pas perpendiculaires alors elles ne se coupent pas}}&1,% 
		{\textbf{Réciproque \textcolor{white}{si les droites sont perpendiculaires alors elles se coupent}}}&2,% 
	} 
	%\noindent
	%\QCM[VF,Alterne,Stretch=1.1, Largeur=0.75cm]{%
	%	{\textbf{Implication} Si \uline{le quadrilatère est un parallélogramme} alors \uline{ses diagonales ont le même milieu}.}&1,%
	%	{\textbf{Contraposée} \textcolor{gray!15}{Si les diagonales n'ont pas le même milieu alors le quadrilatère n'est pas un parallélogramme}}&1,% 
	%	{\textbf{Réciproque} \textcolor{white}{Si les diagonales ont le même milieu alors le quadrilatère est un parallélogramme}}&2,% 
	%}  
\end{exercise} 
\begin{tBox}
	Une implication et sa contraposée ont toujours même valeur de vérité (soit toutes les deux vraies, soit toutes les deux fausses).
\end{tBox}
\begin{definition}
	Un nombre $p\in\mathbb{N}$ est pair (multiple de 2) s'il peut s'écrire sous la forme $p=2k$, avec $k\in\mathbb{N}$.
	
	Un nombre $q\in\mathbb{N}$ est impair, si le reste de la division de $q$ par 2 est 1. c.à.d. il peut s'écrire sous la forme $2k+1$ avec $k\in\mathbb{N}$.
\end{definition}
\begin{theorem} Pour tout entier $n\in\mathbb{N}$ :
	
	\noindent\parbox{0.48\linewidth}{
		Si $n$ est pair alors $n^2$ est pair.
	}\hfill \parbox{0.48\linewidth}{
		Si $n$ est impair alors $n^2$ est impair.
	}
\end{theorem}
\begin{exercise}[au programme] Démontrer les implications.
	
	\noindent\parbox{0.48\linewidth}{
		$n$ est pair. Il existe $k\in\mathbb{N}$ tel que
		\[
		n=2k
		\]
	}\hfill \parbox{0.48\linewidth}{
		$n$ est impair. Il existe $k\in\mathbb{N}$ tel que
		\[
		n=2k+1
		\] 
	}
	\vfill 
	\vfill 
\end{exercise}
\begin{theorem} Pour tout entier $n\in\mathbb{N}$ :
	
	\noindent\parbox{0.48\linewidth}{
		Si $n^2$ est pair alors $n$ est pair.
	}\hfill \parbox{0.48\linewidth}{
		Si $n^2$ est impair alors $n$ est impair.
	} 
\end{theorem}
\noindent\parbox{0.48\linewidth}{ 
	\paragraph{Contraposée} Si
}\hfill \parbox{0.48\linewidth}{ 
	\paragraph{Contraposée} Si
}
\vfill 
%----------------------------------------------------------------------------------------
%	
%----------------------------------------------------------------------------------------